\documentclass[11pt,a4paper]{article}
\usepackage[T1]{fontenc}
\usepackage[utf8]{inputenc}
\usepackage{lmodern}
\usepackage[margin=27mm,headheight=14pt]{geometry}
\usepackage{amsmath,amssymb,amsthm,mathtools}
\usepackage{microtype,booktabs,array,longtable}
\usepackage{xcolor}
\usepackage{enumitem}
\usepackage{fancyhdr}
\usepackage[colorlinks=true,linkcolor=blue!45!black,citecolor=blue!45!black,urlcolor=blue!45!black]{hyperref}
\usepackage[nameinlink,noabbrev]{cleveref}
\hypersetup{pdftitle={Recovering Uniform Factors: Exact Moment Fibres and Logarithmic Instability Near the Fabius--Rvachev Law},pdfauthor={Research draft prepared with ChatGPT},pdfsubject={Inverse problems for infinite convolutions of uniform distributions}}
\pagestyle{fancy}
\fancyhf{}
\fancyhead[L]{\small Recovering uniform factors}
\fancyhead[R]{\small Research draft}
\fancyfoot[C]{\thepage}
\renewcommand{\headrulewidth}{0.3pt}
\setlength{\emergencystretch}{2em}
\setlength{\parskip}{3pt}
\setlist[enumerate]{itemsep=3pt,topsep=4pt}
\newtheorem{theorem}{Theorem}[section]
\newtheorem{proposition}[theorem]{Proposition}
\newtheorem{lemma}[theorem]{Lemma}
\newtheorem{corollary}[theorem]{Corollary}
\theoremstyle{definition}
\newtheorem{definition}[theorem]{Definition}
\newtheorem{example}[theorem]{Example}
\newtheorem{question}{Research question}
\theoremstyle{remark}
\newtheorem{remark}[theorem]{Remark}
% ed. (2026-09-29): unnumbered environment for editorial notes.
\newtheorem*{ednote}{Editorial note (ProveIt, 2026-09-29)}
\DeclareMathOperator{\TV}{TV}
\DeclareMathOperator{\Law}{Law}
\DeclareMathOperator{\Var}{Var}
\DeclareMathOperator{\supp}{supp}
\DeclareMathOperator{\sinc}{sinc}
\DeclareMathOperator{\diag}{diag}
\newcommand{\R}{\mathbb R}
\newcommand{\N}{\mathbb N}
\newcommand{\E}{\mathbb E}
\newcommand{\A}{\mathcal A}
\newcommand{\K}{\mathcal K}
\newcommand{\dd}{\,\mathrm d}
\newcommand{\eps}{\varepsilon}
\newcommand{\norm}[1]{\left\lVert #1\right\rVert}
\newcommand{\abs}[1]{\left\lvert #1\right\rvert}
\newcommand{\up}{\operatorname{up}}
\newcommand{\sort}{\operatorname{sort}_{\downarrow}}
\newcommand{\repoSHA}{74f7f5bdba1aa728b340509701a0fa4e45e8dbf3}

\title{\textbf{Recovering Uniform Factors}\\[4pt]
\large Exact Moment Fibres and Logarithmic Instability\\
Near the Fabius--Rvachev Law}
\author{Research article prepared for Vladimir Reshetnikov\\
\small AI-assisted mathematical draft; proofs and reproducible checks included}
\date{September 2026}

\begin{document}
\maketitle

\begin{abstract}
Let $X_a=\sum_{j\ge1}a_jU_j$, where the $U_j$ are independent uniform random
variables on $[-1,1]$ and $a$ is a decreasing, nonnegative, summable sequence.
The distribution of $X_a$ determines $a$ uniquely; this is an established
identifiability theorem, not a new claim of this article. We investigate how
much information survives finite moments and small errors in the distribution.

First, for every $M$, we construct a nonconstant real-analytic family through
the dyadic Fabius--Rvachev spectrum that preserves its exact support and all
moments through degree $2M+1$. Second, Chebyshev level sets give explicit
pairs $a_n^+,a_n^-$ whose laws have equal moments through degree $2n-1$ but
satisfy
\[
 d_\infty(a_n^+,a_n^-)\ge\frac1{48n^2},\qquad
 \TV(\Law(X_{a_n^+}),\Law(X_{a_n^-}))
 \le 40n(6/7)^{2n-1}.
\]
The laws converge in total variation to the Fabius--Rvachev law. A common
normalization fixes their variance at $1/9$, keeps support inside $[-3,3]$,
and preserves uniform bounds on every derivative of their densities; the
spectral lower bound then becomes $1/(96n^2)$.
Consequently the inverse is not pairwise H\"older on any total-variation
neighborhood of that law. Its local modulus is bounded below by a constant
times $(\log(1/\eps))^{-2}$, and estimation from $N$ independent observations
has a corresponding minimax lower bound of order $(\log N)^{-2}$.
We give complete proofs, positive recovery results under finite-dimensional
separation assumptions, and a reproducible exact-arithmetic verification
suite. The exponent is a lower-bound exponent, not a proved optimal rate.
\end{abstract}

\noindent\textbf{Status.} The quantitative construction is presented as a proposed
research contribution. Its priority has not been established by an exhaustive
literature search. No claim to solve a named, previously published open
conjecture is made. The proofs below are conventional and unrefereed; neither
they nor the accompanying computations constitute a Lean formalization.

\clearpage
\tableofcontents
\clearpage

\section{The inverse question and its relation to ProveIt}
\label{sec:intro}

\subsection{From a forward construction to an inverse problem}

The Fabius-function part of ProveIt develops the Rvachev up-function,
infinite products of sinc functions, geometric-uniform sums, and their
moment identities. Its project documentation explicitly identifies the two
papers of Arias de Reyna as foundational sources
\cite{ProveIt,AriasDefinition,AriasArithmetic}. The natural forward question
is to determine a distribution, its transform, or its moments from the
lengths of its independent uniform summands. Here we reverse that direction:
\emph{how reliably can those lengths be recovered from the distribution?}

This question has three different meanings. With the entire distribution
known exactly, identifiability is already understood. With only finitely
many exact moments, there can be infinitely many different answers. With
an approximately known entire distribution, an inverse can be continuous
and nevertheless fail every power-law error bound. The article separates
these issues rather than using the word ``recoverable'' for all three.

Billey and Swanson prove identifiability for generalized uniform sums in a
larger square-summable parameter space; see their Theorem~3.13 and
Lemma~3.11 \cite{BilleySwanson}. Their Theorem~1.15 gives a compact
homeomorphic parametrization after allowing a Gaussian component in a
suitable standardized space. We do not claim those results as new. The
specific contribution here is an explicit, quantitatively ill-conditioned
family inside the compact-support uniform-convolution class, localized near
a distinguished smooth law already central to ProveIt.

Classical moment inversion is also known to be ill-conditioned. For example,
Gerth, Hofmann, Hofmann, and Kindermann analyze conditional stability and
failure of H\"older stability for the Hausdorff moment problem
\cite{GerthEtAl}. That general theory does not by itself produce the objects
needed here: every parameter must be a genuine multiset of positive uniform
half-lengths, the observable is the full probability law rather than just a
sequence of noisy moments, and the laws must remain uniformly smooth.
Our construction meets those constraints directly.

\subsection{Notation and the main conclusion}

Let $U_1,U_2,\ldots$ be independent and uniform on $[-1,1]$. For $L>0$ set
\[
 \A_L=\left\{a=(a_j)_{j\ge1}:a_1\ge a_2\ge\cdots\ge0,
                 \quad\sum_{j\ge1}a_j\le L\right\}.
\]
Define
\[
 X_a=\sum_{j\ge1}a_jU_j,\qquad \mu_a=\Law(X_a),\qquad
 d_\infty(a,b)=\sup_{j\ge1}|a_j-b_j|.
\]
Zero coefficients are padding, not additional identifiable factors. The sum
converges absolutely for every outcome. We use
\[
 \TV(\mu,\nu)=\sup_{B\text{ Borel}}|\mu(B)-\nu(B)|
             =\frac12\norm{f-g}_1
\]
when densities $f,g$ exist. All logarithms without a subscript are natural.
Write
\[
 a^0=(2^{-1},2^{-2},\ldots),\qquad \mu_0=\mu_{a^0},\qquad f=\up
\]
for the density of $\mu_0$. In this normalization the support is $[-1,1]$,
$f(0)=1$, and $\Var(X_{a^0})=1/9$; these facts are proved in
\cref{sec:fabius}.

\begin{definition}\label{def:K}
Let $\K$ be the subset of $\A_3$ consisting of spectra whose laws have
variance $1/9$ and a density $g\in C_c^\infty(\R)$ satisfying, for every
integer $r\ge0$,
\[
 \norm{g^{(r)}}_\infty\le 2^{r+1}\norm{f^{(r)}}_\infty.
\]
Thus $\K$ imposes a common support bound, an exact variance, and a fixed
bound for each derivative order. It contains $a^0$.
\end{definition}

\begin{theorem}[Quantitative instability near the dyadic law]
\label{thm:main}
For every $n\ge2$ there are explicit spectra $\widetilde a_n^\pm\in\K$
with the following properties:
\begin{align}
 d_\infty(\widetilde a_n^+,\widetilde a_n^-)&\ge\frac1{96n^2},
                                                    \label{eq:main-gap}\\
 \TV(\mu_{\widetilde a_n^+},\mu_{\widetilde a_n^-})
     &\le \eps_n:=\min\{1,40n(6/7)^{2n-1}\},          \label{eq:main-tv}\\
 \TV(\mu_{\widetilde a_n^\pm},\mu_0)
     &\le \sqrt{\frac5{12n}}+\frac{45}{16n}.           \label{eq:main-near}
\end{align}
The two laws have identical moments through degree $2n-1$ and different
moments in degree $2n$.

For every $r>0$ and every $C,\alpha>0$, the estimate
\[
 d_\infty(a,b)\le C\TV(\mu_a,\mu_b)^\alpha
\]
fails for some $a,b\in\K$ with both laws in the total-variation $r$-ball
about $\mu_0$. Moreover, for all sufficiently small $\eps>0$, the local
inverse modulus in that ball is at least
\[
 \frac{1}{11616[\log(100/\eps)]^2}.
\]
\end{theorem}

The proof occupies \cref{sec:chebyshev,sec:fourier,sec:fixed,sec:consequences}.
The important feature is not the numerical constant. It is the contrast
between a polynomial separation of parameters and an exponentially small
separation of full laws, under simultaneous smoothness and variance control.

\subsection{What is established, proposed, and left unresolved}

\begin{center}
\small
\begin{tabular}{@{}p{0.31\textwidth}p{0.62\textwidth}@{}}
\toprule
Statement & Status in this article\\
\midrule
Cumulants and exact identifiability & Established background; re-derived with
attribution to Billey--Swanson.\\
Finite moment fibres through $a^0$ & A proved local construction; not advertised
as a named open-conjecture solution.\\
Chebyshev pairs and explicit TV bound & Proposed quantitative contribution;
proved for every $n$, independently of numerical tests.\\
Fixed-variance, uniformly smooth instability & Consequence of the explicit
construction, with all normalization constants given.\\
Logarithmic statistical obstruction & A lower bound proved by a two-point
argument, not an optimal minimax rate.\\
Optimal inverse modulus and pointwise stability at $a^0$ & Not settled here;
precise questions appear in \cref{sec:questions}.\\
\bottomrule
\end{tabular}
\end{center}

The repository review was targeted, not an exhaustive audit of ProveIt.
The source reference is pinned to the observed commit
\texttt{\repoSHA}; the accompanying source notes identify which documentation
and outside literature were consulted.

\section{Transforms, cumulants, and qualitative recovery}
\label{sec:background}

\subsection{Support and the power-sum encoding}

\begin{proposition}\label{prop:cumulants}
For $a\in\A_L$, put $S(a)=\sum_j a_j$ and $p_k(a)=\sum_j a_j^{2k}$.
Then
\begin{align}
 \supp\mu_a&=[-S(a),S(a)],                            \label{eq:support}\\
 \phi_a(t):=\E e^{itX_a}&=\prod_{j\ge1}\sinc(a_jt),  \label{eq:char}\\
 \kappa_{2k}(X_a)&=\frac{2^{2k}B_{2k}}{2k}p_k(a)
       \quad(k\ge1),                                \label{eq:cumulants}\\
 \kappa_{2k+1}(X_a)&=0\quad(k\ge0).                  \label{eq:odd}
\end{align}
Here $\sinc z=\sin z/z$, with its removable value at zero, and $B_m$
are the Bernoulli numbers with $B_2=1/6$ and $B_4=-1/30$.
For $|t|<\pi/a_1$ when $a_1>0$,
\begin{equation}\label{eq:log-sinc}
 \log\phi_a(t)
 =-\sum_{k\ge1}\frac{\zeta(2k)}{k\pi^{2k}}p_k(a)t^{2k}.
\end{equation}
\end{proposition}

\begin{proof}
The map from the product cube $[-1,1]^{\N}$ to $\R$ given by
$(u_j)\mapsto\sum_j a_ju_j$ is continuous: its tails are bounded uniformly
by $\sum_{j>N}a_j$. Its image is the entire interval in
\eqref{eq:support}, since the constant sequence $u_j=x/S(a)$ realizes any
point of that interval when $S(a)>0$. Product uniform measure gives positive
measure to every nonempty basic open cylinder. Continuity therefore implies
that every point of the image is in the support of its pushforward. The
case $S(a)=0$ is immediate.

Independence gives \eqref{eq:char} first for partial sums and then by
bounded convergence. Euler's product for sine yields
\[
 \log\sinc z=-\sum_{k\ge1}\frac{\zeta(2k)}{k\pi^{2k}}z^{2k},
                  \qquad |z|<\pi.
\]
For $|t|<\pi/a_1$, the double series is absolutely convergent. Indeed,
$p_k(a)\le a_1^{2k-2}p_1(a)$, so after extracting $t^2p_1(a)$ the
remaining series is geometrically dominated. This justifies summation over
$j$ and gives \eqref{eq:log-sinc}.

Alternatively expand
$\log(\sinh z/z)=\sum_{k\ge1}2^{2k}B_{2k}z^{2k}/(2k(2k)!)$
near zero and use additivity of log moment-generating functions. The
coefficient of $z^{2k}/(2k)!$ is \eqref{eq:cumulants}. Symmetry gives
\eqref{eq:odd}. In particular, $\Var(X_a)=p_1(a)/3$.
\end{proof}

Moment and cumulant lists determine each other triangularly. Thus equality
of $p_1,\ldots,p_M$ is equivalent to equality of even cumulants through
order $2M$, and it implies equality of all moments through degree $2M+1$.
The extra odd degree uses symmetry. Since every $B_{2k}$ is nonzero, no
even power sum disappears from the full cumulant sequence.

\subsection{The known uniqueness theorem}

\begin{theorem}[Identifiability; a summable special case of Billey--Swanson]
\label{thm:unique}
If $a,b\in\A_L$ and $\mu_a=\mu_b$, then $a=b$.
\end{theorem}

\begin{proof}
Equal laws have equal cumulants and hence equal $p_k$. For any nonzero
spectrum let $A=a_1>0$. Then
\[
 A^{2k}\le p_k(a)\le A^{2k-2}p_1(a),
\]
so $p_k(a)^{1/(2k)}\to A$. Its multiplicity is finite and equals
\[
 \lim_{k\to\infty}\frac{p_k(a)}{A^{2k}}
 =\sum_j\lim_{k\to\infty}(a_j/A)^{2k}
 =\#\{j:a_j=A\}.
\]
Dominated convergence is justified by the summable majorant $(a_j/A)^2$.
Remove these largest equal entries from both spectra and subtract their
contribution from every power sum. Repeat. Every positive coefficient at a
fixed rank is reached after finitely many such removals. If a remaining
power-sum sequence vanishes, all its coefficients vanish. This recovers the
whole padded decreasing sequence.
\end{proof}

The proof is included to specify exactly what the inverse means.
It does not turn the already known theorem into a new result. Nor is its
limit procedure a robust numerical algorithm: taking high cumulants and
successive deflations can amplify observational errors severely.

\subsection{Continuity without a useful power-law modulus}

\begin{proposition}\label{prop:compact}
The space $\A_L$ is compact in $\ell^2$. Its map to laws is continuous in
the quadratic Wasserstein distance $W_2$, with
\begin{equation}\label{eq:forward-w2}
 W_2(\mu_a,\mu_b)\le\frac{\norm{a-b}_2}{\sqrt3}.
\end{equation}
It is a homeomorphism onto its image in $W_2$ and, equivalently on this
bounded-support image, in the weak topology.
\end{proposition}

\begin{proof}
Monotonicity implies $a_j\le L/j$. Consequently
$\sum_{j>N}a_j^2\le L^2\sum_{j>N}j^{-2}\le L^2/N$, uniformly in
$a\in\A_L$. A coordinatewise diagonal subsequence therefore converges in
$\ell^2$. Its limit is decreasing and nonnegative, and has sum at most $L$
by Fatou's lemma. This proves compactness.

Couple both random series with the same $U_j$. Independence and centering
give $\E|X_a-X_b|^2=\norm{a-b}_2^2/3$, which proves
\eqref{eq:forward-w2}. By \cref{thm:unique}, the continuous map is
injective. A continuous injection from a compact space into a Hausdorff
space is a homeomorphism onto its image. Finally, all supports are in
$[-L,L]$, where weak convergence implies convergence of second moments
and hence $W_2$ convergence.
\end{proof}

In particular an inverse modulus tending to zero exists. The main theorem
concerns its \emph{rate}. There is no contradiction between a compact
homeomorphism and a failure of every H\"older estimate. We make no claim that
the forward map is globally continuous in total variation: approaching a
point mass is an elementary obstruction to such a statement.

\begin{proposition}[A largest-factor certificate]\label{prop:largest}
If $a\in\A_L$, $L>0$, and $k\ge1$, then
\begin{equation}\label{eq:largest}
 \left(\frac{p_k(a)}L\right)^{1/(2k-1)}
       \le a_1\le p_k(a)^{1/(2k)}.
\end{equation}
The length of this interval is at most $L/(2k)$.
\end{proposition}

\begin{proof}
The upper estimate uses $a_1^{2k}\le p_k$; the lower uses
$p_k\le a_1^{2k-1}\sum_j a_j\le La_1^{2k-1}$.
Put $q=p_k/L^{2k}\in[0,1]$. The interval length divided by $L$ is
$q^{1/(2k)}-q^{1/(2k-1)}$. Its maximum on $[0,1]$ is
\[
 \frac1{2k}\left(\frac{2k-1}{2k}\right)^{2k-1}\le\frac1{2k},
\]
by differentiation, including the endpoint values zero.
\end{proof}

If a certified measurement satisfies $|\widehat p_k-p_k|\le\delta$, the
same proof gives the enclosing interval
\[
 \left[
  \left(\frac{\max\{0,\widehat p_k-\delta\}}{L}\right)^{1/(2k-1)},
  (\widehat p_k+\delta)^{1/(2k)}
 \right]\cap[0,L].
\]
This is a statement about a \emph{certified power-sum error}. It does not
supply such an error from raw empirical moments, and it does not establish
a rate for recovery of the entire spectrum.

\section{The smooth dyadic reference law}
\label{sec:fabius}

\begin{proposition}\label{prop:base}
The random variable $X_0=\sum_{j\ge1}2^{-j}U_j$ has an even, unimodal
density $f\in C_c^\infty(\R)$ with
\[
 \supp f=[-1,1],\qquad f(0)=1,\qquad
 \norm{f'}_1=2,\qquad \Var(X_0)=\frac19.
\]
This is the usual Rvachev up-density in the normalization of
\cite{AriasDefinition}.
\end{proposition}

\begin{proof}
Support and variance follow from \cref{prop:cumulants} and
$\sum_{j\ge1}4^{-j}=1/3$. For any fixed $m$, retaining the first $m$
factors in the characteristic product gives
\[
 |\phi_0(t)|\le \prod_{j=1}^m\min\{1,2^j/|t|\}
               \le C_m(1+|t|)^{-m}.
\]
Thus $t^r\phi_0(t)$ is integrable for every $r$. Fourier inversion and
differentiation under the integral give a $C^\infty$ density, necessarily
compactly supported.

For completeness, even unimodality is preserved by convolution with a
centered uniform density. If $g$ is even and nonincreasing on $[0,\infty)$,
then the derivative of $(2a)^{-1}\int_{x-a}^{x+a}g(u)\dd u$ is
$(g(x+a)-g(x-a))/(2a)\le0$ for $x\ge0$, wherever the derivative exists.
Induction applies to finite convolutions. Their characteristic functions
are dominated, from the second convolution onward, by an integrable
multiple of $(1+|t|)^{-2}$. Fourier inversion therefore gives uniform
convergence of their densities to $f$. Even unimodality passes to the limit.

Write $X_0=U_1/2+Y$, where $Y$ is supported in $[-1/2,1/2]$ and has no
atoms, because it itself has an independent nondegenerate uniform summand.
The first summand has density $1$ on $[-1/2,1/2]$. Its convolution with
$Y$ gives $f(0)=1$. Smoothness, evenness, monotonicity, and compact support
then yield $\norm{f'}_1=2f(0)=2$.
\end{proof}

\begin{lemma}[Translation and smoothing estimates]\label{lem:smoothing}
If $Y$ is independent of $X_0$ and integrable, then
\begin{equation}\label{eq:translation}
 \TV(\Law(X_0+Y),\mu_0)\le \E|Y|.
\end{equation}
If $Y$ is centered and square-integrable, the right side is at most
$\sqrt{\Var Y}$. For an arbitrary probability measure $\nu$,
\begin{equation}\label{eq:smooth-bound}
 \norm{(f*\nu)^{(r)}}_\infty\le\norm{f^{(r)}}_\infty
            \qquad(r\ge0).
\end{equation}
\end{lemma}

\begin{proof}
The fundamental theorem of calculus and Fubini give
$\norm{f(\cdot-y)-f}_1\le |y|\norm{f'}_1=2|y|$.
Average this inequality over $Y$ and divide by two. Cauchy--Schwarz gives
the variance bound. Differentiation under convolution with a probability
measure proves \eqref{eq:smooth-bound}.
\end{proof}

The role of the common Fabius factor is now clear. It makes every
constructed law smooth without demanding that the perturbing finite sum
be smooth to all orders, and it provides explicit total-variation control
of convergence back to the reference law.

\section{Exact finite-moment fibres with fixed support}
\label{sec:fibres}

\subsection{A rank argument that preserves positivity}

\begin{lemma}[Generalized Vandermonde rank]\label{lem:vander}
For distinct positive $t_1,\ldots,t_d$ and distinct nonnegative integers
$e_1<\cdots<e_d$, the matrix $(t_j^{e_i})_{i,j=1}^d$ is invertible.
\end{lemma}

\begin{proof}
A nonzero linear combination of $d$ distinct monomials has at most $d-1$
distinct positive zeros. Prove this by induction on $d$: divide by the
smallest power of $t$ without changing positive zeros, differentiate, and
apply Rolle's theorem. The derivative has at most $d-1$ nonzero monomials,
so has at most $d-2$ positive zeros. If a linear combination of the rows
vanished at all $d$ nodes, this zero bound would be contradicted.
\end{proof}

\begin{theorem}[A curve of laws invisible to any prescribed finite moment list]
\label{thm:fibres}
For every integer $M\ge1$, there are $\delta>0$ and a nonconstant
real-analytic curve $\theta\mapsto a(\theta)\in\A_1$, $|\theta|<\delta$,
such that $a(0)=a^0$ and:
\begin{enumerate}[label=(\roman*)]
\item only the first $M+2$ coefficients change;
\item $\supp\mu_{a(\theta)}=[-1,1]$ exactly;
\item all moments through degree $2M+1$ are independent of $\theta$;
\item the cumulant of order $2M+2$ has a nonzero derivative at zero after
      a suitable local parametrization;
\item the laws have $C_c^\infty$ densities and tend to $\mu_0$ in total
      variation as $\theta\to0$.
\end{enumerate}
\end{theorem}

\begin{proof}
Set $r=M+2$ and let $t=(t_1,\ldots,t_r)$ denote the changing half-lengths.
Consider the analytic map
\[
 F(t)=\left(\sum_{j=1}^r t_j,\ \sum_{j=1}^r t_j^2,
                 \ \sum_{j=1}^r t_j^4,\ldots,
                    \sum_{j=1}^r t_j^{2M}\right).
\]
At $t^0=(2^{-1},\ldots,2^{-r})$, its Jacobian has full row rank $M+1$.
After nonzero row factors are removed, any square maximal minor is a
matrix from \cref{lem:vander}, with exponents
$0,1,3,\ldots,2M-1$.

The analytic implicit function theorem gives a one-dimensional analytic
level manifold through $t^0$. Positivity, strict decreasing order, and the
gap above the unchanged coefficient $2^{-r-1}$ are open conditions, so a
sufficiently short curve in this manifold remains a valid spectrum.
Constancy of its first coordinate preserves $\sum_j a_j=1$, proving (ii).
Constancy of the other coordinates preserves $p_1,\ldots,p_M$ and hence
(iii).

Append $G(t)=\sum_j t_j^{2M+2}$ to $F$. The resulting square Jacobian is
invertible by the same lemma, now with exponents
$0,1,3,\ldots,2M+1$. Thus $\dd G$ does not vanish on the one-dimensional
kernel of $\dd F$. We may use $G-G(t^0)$ itself as a local parameter on
the level manifold. Since $B_{2M+2}\ne0$, \eqref{eq:cumulants} proves (iv).
In particular, the laws are genuinely different; this is not a
reordering of identical factors.

The unchanged tail is distributed as $2^{-r}X_0$. It has a smooth compact
support density $g(x)=2^r f(2^rx)$. All laws on the curve are convolutions
with this common density. Couple the finite changing prefixes with the
same $U_1,\ldots,U_r$. Their difference tends to zero in $L^1$, and
$\norm{g(\cdot-y)-g}_1\le |y|\norm{g'}_1$ proves total-variation
convergence. Smoothness follows by convolution, completing (v).
\end{proof}

More generally, changing $r$ distinct positive coefficients while fixing
the support radius and the first $M$ even power sums gives a local
manifold of dimension $r-M-1$ whenever $r\ge M+1$. In particular, $M+2$
is the smallest number of distinct changing coefficients for which this
regular local argument gives a nonconstant curve. This is a local rank
statement, not a global classification of all moment fibres.

\subsection{A completely explicit fixed-variance example}

\begin{example}\label{ex:curve}
For $M=1$, leave $a_j=2^{-j}$ for $j\ge4$ and set
\begin{align*}
 a_1(\theta)&=(7+5\cos\theta-\sqrt3\sin\theta)/24,\\
 a_2(\theta)&=(7-\cos\theta+3\sqrt3\sin\theta)/24,\\
 a_3(\theta)&=(7-4\cos\theta-2\sqrt3\sin\theta)/24.
\end{align*}
For sufficiently small $|\theta|$ these entries are positive, strictly
decreasing, and greater than $1/16$. Direct calculation gives
\[
 \sum_{j=1}^3a_j(\theta)=\frac78,\qquad
 \sum_{j=1}^3a_j(\theta)^2=\frac{21}{64}.
\]
The full support is exactly $[-1,1]$ and the variance remains $1/9$.
Nevertheless
\[
 \left.\frac{\dd}{\dd\theta}\kappa_4(X_{a(\theta)})\right|_{\theta=0}
       =\frac{7\sqrt3}{3840}\ne0.
\]
Thus even exact support and exact variance do not identify the Fabius
spectrum. The verification script checks these identities symbolically.
\end{example}

The finite-fibre theorem and the quantitative theorem below answer
\emph{different} questions. The former fixes the exact support endpoints
and agrees with $\mu_0$ on a chosen finite moment list. The latter compares
two laws that are exponentially close to one another in total variation;
it imposes a common support bound, not exact support endpoints $\pm1$.
Neither statement should be silently substituted for the other.

\section{Chebyshev pairs of genuine uniform-factor spectra}
\label{sec:chebyshev}

\subsection{Two root sets with exactly matching initial power sums}

Let $T_n$ be the Chebyshev polynomial characterized by
$T_n(\cos\theta)=\cos(n\theta)$. For $n\ge2$, list in decreasing order
\[
 y_0^+>\cdots>y_{n-1}^+,\qquad
 y_0^->\cdots>y_{n-1}^-
\]
the roots of $T_n(y)=1/2$ and $T_n(y)=-1/2$, respectively. Every root is
simple and in $(-1,1)$. On the $j$th interval of monotonicity in angular
coordinates, the two angles are
\begin{equation}\label{eq:angles}
 \theta_j^\pm=\frac{j\pi+\alpha_j^\pm}{n},\qquad
 \alpha_j^+=
 \begin{cases}\pi/3&j\text{ even},\\2\pi/3&j\text{ odd},\end{cases}
 \qquad \alpha_j^-=\pi-\alpha_j^+.
\end{equation}
Then $y_j^\pm=\cos\theta_j^\pm$.

\begin{lemma}\label{lem:cheby-powers}
For $1\le k<n$,
\begin{equation}\label{eq:equal-powers}
 \sum_{j=0}^{n-1}(3+y_j^+)^k
 =\sum_{j=0}^{n-1}(3+y_j^-)^k.
\end{equation}
At the first unmatched order,
\begin{equation}\label{eq:leading-defect}
 \sum_{j=0}^{n-1}(3+y_j^+)^n
 -\sum_{j=0}^{n-1}(3+y_j^-)^n=n2^{1-n}.
\end{equation}
Both first power sums in \eqref{eq:equal-powers} are $3n$.
\end{lemma}

\begin{proof}
The monic polynomials
\[
 P_\pm(y)=2^{1-n}(T_n(y)\mp1/2)
\]
differ only in their constant coefficient, with
$[y^0](P_+-P_-)=-2^{1-n}$. Newton's identities say that their root power
sums agree in degrees $1,\ldots,n-1$; at degree $n$ the difference is
$-n$ times this constant-coefficient difference, namely $n2^{1-n}$.
The binomial theorem transports both assertions from $y$ to $3+y$:
all lower-degree differences cancel. The coefficient of $y^{n-1}$ in
$T_n$ vanishes, so the root sum is zero, giving $3n$ after translation.
\end{proof}

The translated roots lie in $(2,4)$. Taking their positive square roots
therefore produces actual uniform half-lengths, not a signed measure or
formal perturbation. This positivity is essential to the inverse problem.

\subsection{Insertion into an empty dyadic scale interval}

Define
\begin{equation}\label{eq:eta}
 k_n=\lceil\log_2 n\rceil,\qquad
 \eta_n=2^{-k_n-1},\qquad
 \frac1{4n}\le\eta_n\le\frac1{2n}.
\end{equation}
Let all uniforms in the following expression be independent, and independent
of $X_0$:
\begin{equation}\label{eq:S}
 S_n^\pm=
 \sum_{j=0}^{n-1}\sqrt{3+y_j^\pm}\,U_j
          +\sum_{\ell=1}^{2n}V_\ell,
 \qquad X_n^\pm=X_0+\eta_n S_n^\pm.
\end{equation}
Let $a_n^\pm$ be the decreasing rearrangement of the multiset
\begin{equation}\label{eq:spectra}
 \{2^{-j}:j\ge1\}
 \ \cup\ \{\eta_n\sqrt{3+y_j^\pm}:0\le j<n\}
 \ \cup\ \{\eta_n\text{ repeated }2n\text{ times}\}.
\end{equation}
The last $2n$ identical factors are not cosmetic. They suppress the
large-frequency part of the Fourier transform, where moment matching
alone would not control total variation.

\begin{proposition}\label{prop:pairs}
The spectra in \eqref{eq:spectra} belong to $\A_3$ and satisfy
\begin{align}
 \Var(X_n^\pm)&=\frac19+\frac{5n\eta_n^2}{3},             \label{eq:variance}\\
 \TV(\Law(X_n^\pm),\mu_0)&\le\sqrt{5n/3}\,\eta_n
                                    \le\sqrt{5/(12n)}. \label{eq:base-close}
\end{align}
Their densities $f_n^\pm$ are smooth, supported in $[-3,3]$, and satisfy
$\norm{(f_n^\pm)^{(r)}}_\infty\le\norm{f^{(r)}}_\infty$ for all $r$.
The laws agree in moments through degree $2n-1$, whereas
\begin{equation}\label{eq:cumulant-defect}
 \kappa_{2n}(X_n^+)-\kappa_{2n}(X_n^-)
         =2^n B_{2n}\eta_n^{2n}\ne0.
\end{equation}
\end{proposition}

\begin{proof}
Each of the $n$ variable half-lengths before scaling is less than $2$,
so the added sum of half-lengths is at most $4n\eta_n\le2$. The base
contributes $1$. Their total squared half-length before scaling is
$3n+2n=5n$ by \cref{lem:cheby-powers}, proving \eqref{eq:variance}.
Apply \cref{lem:smoothing} with $Y=\eta_nS_n^\pm$ to obtain
\eqref{eq:base-close} and all smoothness bounds.

The first $n-1$ even power sums agree by \eqref{eq:equal-powers}; the base
and the repeated factors are common. Hence all moments through degree
$2n-1$ agree. Multiplying \eqref{eq:leading-defect} by
$\eta_n^{2n}2^{2n}B_{2n}/(2n)$ gives \eqref{eq:cumulant-defect}.
The first different cumulant is also the first different moment, since
the moment--cumulant relation is triangular.
\end{proof}

\subsection{A middle coordinate gives a quadratic, not cubic, separation}

\begin{lemma}[Ordered spectral separation]\label{lem:gap}
For all $n\ge2$,
\begin{equation}\label{eq:gap}
 d_\infty(a_n^+,a_n^-)\ge\frac1{48n^2}.
\end{equation}
More precisely, put $j_n=\lfloor(n-1)/2\rfloor$. A witnessing coordinate
is $k_n+1+j_n$, whose absolute difference is
\begin{equation}\label{eq:actual-gap}
 \Delta_n=\eta_n\left|
    \sqrt{3+y_{j_n}^+}-\sqrt{3+y_{j_n}^-}\right|.
\end{equation}
\end{lemma}

\begin{proof}
All inserted variable half-lengths lie strictly between $\eta_n$ and
$2\eta_n$. There are exactly $k_n$ base dyadic factors at least
$2\eta_n$, and no base factors strictly between these endpoints.
All repeated factors equal the lower endpoint. Therefore the ordered
positions $k_n+1,\ldots,k_n+n$ consist exactly of the inserted variable
factors, in their displayed order. Sorting does not erase the difference.

The two angles at position $j_n$ have distance $\pi/(3n)$ and midpoint
\[
 m_n=\frac{(j_n+1/2)\pi}{n}\in[\pi/4,\pi/2].
\]
Consequently
\[
 |y_{j_n}^+-y_{j_n}^-|
 =2\sin m_n\sin\frac{\pi}{6n}
 \ge 2\cdot\frac12\cdot\frac1{3n}
 =\frac1{3n},
\]
where $\sin x\ge2x/\pi$ for $0\le x\le\pi/2$. Rationalizing the square
roots in \eqref{eq:actual-gap}, their denominator is at most $4$.
Thus $\Delta_n\ge\eta_n/(12n)\ge1/(48n^2)$.
\end{proof}

Using only the largest inserted factor would yield a weaker $n^{-3}$
bound, because cosine is flat at its endpoint. The middle coordinate
avoids that loss. Along $n=2^m$, elementary expansion of
\eqref{eq:actual-gap} even gives
\[
 \Delta_n\sim\frac{\pi}{12\sqrt3}\,n^{-2}.
\]
This is an asymptotic for the displayed witnessing coordinate, not an
assertion that the full $d_\infty$ distance equals that coordinate.

\section{From moment matching to exponential total-variation closeness}
\label{sec:fourier}

The next theorem is the analytic core. It does not use an asymptotic
approximation to the density. In particular, its very small bounds are
not obtained by subtracting nearly equal floating-point density values.

\begin{theorem}[A uniform Fourier certificate]\label{thm:tv}
For $n\ge2$ and $S_n^\pm$ as in \eqref{eq:S},
\begin{equation}\label{eq:TV-S}
 \TV(\Law(S_n^+),\Law(S_n^-))
       \le40n(6/7)^{2n-1}.
\end{equation}
The same bound holds for $X_n^\pm$.
\end{theorem}

\begin{proof}
Put $\rho=6/7$, $q=4/\pi^2<1/2$, and let
\[
 F_\pm(t)=\left(\prod_{j=0}^{n-1}
       \sinc\bigl(\sqrt{3+y_j^\pm}\,t\bigr)\right)\sinc(t)^{2n},
 \qquad D(t)=F_+(t)-F_-(t).
\]
We estimate the Fourier $H^1$ norm, treating small and large frequencies
separately.

\paragraph{Small frequencies: $|t|\le1$.}
Every sinc argument has magnitude at most $2<\pi$, so all factors are
positive. Let $L_\pm(t)=\log F_\pm(t)\le0$ and
$H=L_+-L_-$. The common factors cancel from $H$.
For any $k\ge1$ the difference between the variable power sums has
absolute value at most $2n4^k$; for $k<n$ it is zero. Hence
\begin{align*}
 |H(t)|
 &\le4n\sum_{k\ge n}\frac{q^k|t|^{2k}}{k}
 \le\frac4{1-q}q^n|t|^{2n}
 \le8q^n|t|^{2n},\\
 |H'(t)|
 &\le8n\sum_{k\ge n}q^k|t|^{2k-1}
 \le16nq^n|t|^{2n-1}.
\end{align*}
We used $\zeta(2k)<2$ and the convergent log-sinc expansion; termwise
differentiation is justified by geometric domination on this interval.

For $|s|\le2$, $\sinc(s)\ge1/3$. Indeed sinc decreases on $[0,2]$,
and the alternating sine series gives $\sin2\ge2-2^3/6=2/3$.
Also $|\sinc'(s)|\le\E|U|=1/2$, by differentiating the characteristic
function of $U\sim\mathrm{Unif}[-1,1]$. A variable factor therefore
contributes at most $3$ to $|L_\pm'|$, and a common unit factor at most
$3/2$. Thus $|L_\pm'|\le6n$.

The identity
\[
 D=H\int_0^1 e^{(1-u)L_-+uL_+}\dd u
\]
is useful because the exponential is at most $1$, avoiding any unstable
division by $F_\pm$. Its derivative gives
\begin{equation}\label{eq:low-bounds}
 |D(t)|\le8q^n|t|^{2n},\qquad
 |D'(t)|\le64nq^n|t|^{2n-1}.
\end{equation}
Integrating the squares, even the coarse estimate by interval length gives
\begin{equation}\label{eq:low-int}
 \int_{|t|\le1}(|D(t)|^2+|D'(t)|^2)\dd t
       \le8320n^2q^{2n}.
\end{equation}

\paragraph{Large frequencies: $|t|\ge1$.}
Let $h(t)=|\sinc t|\le1$. All the other sinc factors have modulus at
most $1$. The product rule, used without dividing by any factor, gives
\[
 |F_\pm(t)|\le h(t)^{2n},\qquad
 |F_\pm'(t)|\le n h(t)^{2n}+n h(t)^{2n-1}
                     \le2n h(t)^{2n-1}.
\]
Therefore
\begin{equation}\label{eq:high-bounds}
 |D(t)|^2+|D'(t)|^2\le20n^2h(t)^{4n-2}.
\end{equation}
On $1\le |t|\le2$, sinc is positive in magnitude and bounded by
$\sin1<6/7$. For an exact elementary bound, use
$\sin1\le1-1/6+1/120=101/120<6/7$.
On $|t|\ge2$, use $h(t)\le1/|t|$. It follows that
\begin{align*}
 \int_{|t|\ge1}h(t)^{4n-2}\dd t
 &\le2\rho^{4n-2}
             +\frac{2^{-(4n-4)}}{4n-3}
 \le4\rho^{4n-2}.
\end{align*}
The last inequality holds for $n\ge2$, for example by comparing
$(1/2)^{4n-2}$ to $\rho^{4n-2}$. Combining this with
\eqref{eq:high-bounds} gives
\begin{equation}\label{eq:high-int}
 \int_{|t|\ge1}(|D(t)|^2+|D'(t)|^2)\dd t
       \le80n^2\rho^{4n-2}.
\end{equation}
Since $q^{2n}\le\rho^{4n-2}$,
\begin{equation}\label{eq:H1}
 \int_\R(|D|^2+|D'|^2)\dd t
       \le8400n^2\rho^{4n-2}.
\end{equation}

\paragraph{Conversion to total variation.}
Let $g$ be the difference of the densities of $S_n^+$ and $S_n^-$. These
finite sums have bounded compactly supported densities, so $g,xg\in L^2$.
For the Fourier convention $\widehat g(t)=\int e^{itx}g(x)\dd x=D(t)$,
Cauchy--Schwarz and Plancherel yield
\begin{align*}
 \norm{g}_1
 &\le\left(\int_\R\frac{\dd x}{1+x^2}\right)^{1/2}
       \left(\int_\R(1+x^2)|g(x)|^2\dd x\right)^{1/2}\\
 &=\frac1{\sqrt2}
       \left(\int_\R(|D(t)|^2+|D'(t)|^2)\dd t\right)^{1/2}.
\end{align*}
Dividing by two and using \eqref{eq:H1} proves
\[
 \TV(\Law(S_n^+),\Law(S_n^-))
 \le\frac{\sqrt{8400}}{2\sqrt2}\,n\rho^{2n-1}
 <40n\rho^{2n-1}.
\]
Total variation is invariant under a common nonzero dilation and contracts
under convolution with a common probability measure. Apply these facts
first with dilation $\eta_n$ and then convolution by $\mu_0$ to obtain
the assertion for $X_n^\pm$.
\end{proof}

\begin{remark}[Why both frequency ranges are necessary]
Equal low-order moments alone control only a neighborhood of frequency
zero. A proof that stopped with \eqref{eq:low-bounds} would not establish
closeness of densities in $L^1$. The common $2n$ unit factors create the
large-frequency estimate \eqref{eq:high-int}. Keeping a derivative of the
Fourier difference is also essential to the particular weighted
Cauchy--Schwarz argument used here. These two features explain the
architecture of the construction.
\end{remark}

\section{Fixing the variance without sacrificing smoothness}
\label{sec:fixed}

Define the common scale
\begin{equation}\label{eq:cn}
 c_n=(1+15n\eta_n^2)^{-1/2},\qquad
 \widetilde X_n^\pm=c_nX_n^\pm,\qquad
 \widetilde a_n^\pm=c_na_n^\pm.
\end{equation}

\begin{proposition}\label{prop:normalized}
The spectra $\widetilde a_n^\pm$ belong to $\K$ and satisfy all the
quantitative assertions of \cref{thm:main} preceding its non-H\"older
conclusion. Their first unmatched cumulant is
\[
 \kappa_{2n}(\widetilde X_n^+)-\kappa_{2n}(\widetilde X_n^-)
        =c_n^{2n}2^n B_{2n}\eta_n^{2n}\ne0.
\]
\end{proposition}

\begin{proof}
Equation \eqref{eq:variance} shows that the common variance becomes $1/9$.
Since $15n\eta_n^2\le15/(4n)\le15/8$, we have $1/2<c_n\le1$.
Supports remain in $[-3,3]$, and the scaled densities satisfy
\[
 \norm{(\widetilde f_n^\pm)^{(r)}}_\infty
 =c_n^{-r-1}\norm{(f_n^\pm)^{(r)}}_\infty
 \le2^{r+1}\norm{f^{(r)}}_\infty.
\]
The spectral distance is multiplied by $c_n$, proving
\eqref{eq:main-gap}. Total variation between the pair is unchanged,
proving \eqref{eq:main-tv}. Equal moments remain equal under common
scaling, and cumulants scale homogeneously, giving the displayed defect.

It remains to control the distance to the original reference law rather
than to its dilation. Write $f_c(x)=c^{-1}f(x/c)$. Direct differentiation
and a change of variables give
\[
 \norm{\partial_c f_c}_1
 =c^{-1}\int_\R|f(u)+uf'(u)|\dd u
 \le \frac3c,
\]
because $\norm{f}_1=1$, $|u|\le1$ on the support, and $\norm{f'}_1=2$.
Therefore
\[
 \TV(\Law(cX_0),\mu_0)\le\frac32\log(1/c),\qquad0<c\le1.
\]
Using common-dilation invariance in the first term of the triangle
inequality,
\begin{align*}
 \TV(\Law(c_nX_n^\pm),\mu_0)
 &\le\TV(\Law(X_n^\pm),\mu_0)
                +\frac32\log(1/c_n)\\
 &\le\sqrt{\frac5{12n}}
                +\frac34\log(1+15n\eta_n^2)\\
 &\le\sqrt{\frac5{12n}}+\frac{45}{16n}.
\end{align*}
This proves \eqref{eq:main-near}.
\end{proof}

Thus instability is not explained by increasing variance, expanding
support, or loss of smoothness. All three potential explanations have
been removed. The theorem does allow an increasing number of small
factors in one dyadic scale interval; it is this growing unresolved block
that the inverse must disentangle.

\section{Inverse moduli and statistical lower bounds}
\label{sec:consequences}

\subsection{The exact local statement}

For $r>0$ define
\begin{equation}\label{eq:modulus}
 \omega_r(\eps)=\sup\left\{d_\infty(a,b):
 \begin{array}{l}
 a,b\in\K,\quad\TV(\mu_a,\mu_0)<r,\quad\TV(\mu_b,\mu_0)<r,\\
 \TV(\mu_a,\mu_b)\le\eps
 \end{array}\right\}.
\end{equation}

\begin{corollary}\label{cor:modulus}
For every $r>0$, for all sufficiently small $\eps>0$,
\begin{equation}\label{eq:modulus-lower}
 \omega_r(\eps)\ge\frac1{11616[\log(100/\eps)]^2}.
\end{equation}
In particular no pairwise H\"older inverse estimate holds on the indicated
neighborhood, for any positive exponent.
\end{corollary}

\begin{proof}
The non-H\"older assertion already follows by substituting the pairs:
$1/(96n^2)$ eventually exceeds
$C[40n(6/7)^{2n-1}]^\alpha$, whereas both laws approach $\mu_0$.
For the explicit modulus take
\[
 n=\left\lceil10\log(100/\eps)\right\rceil,\qquad0<\eps\le1.
\]
The elementary bound $2\log(7/6)>0.3$ implies
\[
 (6/7)^{2n-1}\le(7/6)e^{-0.3n}
                   \le(7/6)(\eps/100)^3.
\]
Also $n\le51/\eps$, since $\log(100/\eps)\le5/\eps$ on this range.
It follows that $40n(6/7)^{2n-1}\le0.00238\eps^2\le\eps$.
For sufficiently small $\eps$, \eqref{eq:main-near} puts both laws in
the required $r$-ball. Finally
$n\le11\log(100/\eps)$, so \eqref{eq:main-gap} gives
\eqref{eq:modulus-lower}.
\end{proof}

\begin{remark}[Pairwise local versus pointwise at the reference law]
\label{rem:pointwise}
The conclusion rules out a bound that must hold for \emph{every pair}
in a neighborhood. It does not rule out a pointwise estimate
\[
 d_\infty(a,a^0)\le C\TV(\mu_a,\mu_0)^\alpha.
\]
Our construction makes $\mu_{\widetilde a_n^+}$ and
$\mu_{\widetilde a_n^-}$ exponentially close to each other, but the
proved bound on their individual distances to $\mu_0$ is only
polynomial. Those are different comparisons.
\end{remark}

\subsection{A lower bound for recovery from samples}

An estimator based on $N$ observations is a measurable function of
$(Z_1,\ldots,Z_N)$ with values in the space of bounded real sequences.
The loss is $d_\infty(\widehat a,a)$. We do not require the estimator's
output to be a valid uniform-factor spectrum, so the lower bound also
applies to constrained estimators.

\begin{lemma}[Two-point testing bound]\label{lem:two-point}
If parameters $a,b$ have distance $d$ and observation laws $P,Q$, then
for every estimator,
\[
 \max\{\E_Pd_\infty(\widehat a,a),\E_Qd_\infty(\widehat a,b)\}
       \ge\frac d4(1-\TV(P,Q)).
\]
For product laws, $\TV(P^{\otimes N},Q^{\otimes N})\le N\TV(P,Q)$.
\end{lemma}

\begin{proof}
Choose the parameter closer to the estimator's output as a binary test.
On a testing error the estimation loss is at least $d/2$, by the triangle
inequality. The sum of the two testing error probabilities is at least
$1-\TV(P,Q)$, as follows directly by writing the test's acceptance set
in the definition of total variation. The larger risk is at least half
their sum, proving the first assertion. For the second, expand the
difference of product measures as a telescoping sum and use that taking
a product with a probability measure does not change total variation.
\end{proof}

\begin{theorem}[Logarithmic minimax obstruction]\label{thm:statistics}
Let $N\ge1$ and put $n_N=\lceil10\log(100N)\rceil$. Then
\begin{equation}\label{eq:risk}
 \inf_{\widehat a}\ \sup_{a\in\K}
   \E_{\mu_a^{\otimes N}}d_\infty(\widehat a,a)
      \ge\frac1{768n_N^2}.
\end{equation}
For every fixed $r>0$, the same lower bound holds for all sufficiently
large $N$ when the supremum is restricted to
$\{a\in\K:\TV(\mu_a,\mu_0)<r\}$.
Consequently no uniform bound of order $N^{-\beta}$ with $\beta>0$ is
possible on any such local class.
\end{theorem}

\begin{proof}
Use the two normalized spectra of degree $n_N$. The estimate from the
previous proof gives
\[
 N\cdot40n_N(6/7)^{2n_N-1}
 \le\frac{(40)(7/6)(51)}{10^6N}<\frac12,
\]
where $n_N\le51N$. Thus the product observation laws have total
variation at most $1/2$. Substitute their spectral gap $1/(96n_N^2)$
into \cref{lem:two-point}, obtaining $1/(768n_N^2)$.
The local assertion follows from \eqref{eq:main-near}. Since
$n_N=O(\log N)$, the stated exclusion of algebraic rates follows.
\end{proof}

This is an information-theoretic obstruction. It applies before any
restriction on running time, numerical precision, or model architecture.
It is also only a lower bound: the article does not construct an estimator
attaining the logarithmic rate, nor prove that its exponent is optimal.

\section{Where stable recovery does become possible}
\label{sec:positive}

The negative theorem should not be read as saying that every finite
uniform-convolution problem is equally difficult. Its bad pairs use an
increasing number of unknown, closely spaced factors. Fixing their number
and separating their squared lengths changes the situation.

\subsection{Known tails and a finite set of unknown factors}

\begin{proposition}[Finite-factor analytic inversion]\label{prop:finite}
Suppose the entire tail is known and the unknown part consists of exactly
$r$ distinct positive half-lengths $a_1,\ldots,a_r$. After subtracting
the tail contributions, the first $r$ even power sums determine the
unordered unknown multiset exactly. The inverse is locally real-analytic
at every configuration with distinct squared half-lengths.
\end{proposition}

\begin{proof}
Write $x_j=a_j^2$ and $s_k=\sum_{j=1}^r x_j^k$. Newton's recursion
\[
 e_0=1,\qquad
 k e_k=\sum_{i=1}^k(-1)^{i-1}e_{k-i}s_i\quad(1\le k\le r)
\]
constructs the coefficients of
\[
 P(z)=\prod_{j=1}^r(z-x_j)
     =z^r-e_1z^{r-1}+e_2z^{r-2}-\cdots+(-1)^re_r.
\]
Its roots give the multiset. The Jacobian of
$(x_1,\ldots,x_r)\mapsto(s_1,\ldots,s_r)$ has determinant
\[
 r!\prod_{1\le i<j\le r}(x_j-x_i),
\]
up to the harmless sign determined by ordering. It is nonzero at distinct
nodes. The inverse function theorem, followed by positive square roots,
gives the local analytic inverse.
\end{proof}

\begin{proposition}[An explicit root-location certificate]\label{prop:rouche}
Suppose the squared half-lengths satisfy
$m^2\le x_i\le A^2$ and $|x_i-x_j|\ge\gamma>0$ for $i\ne j$.
Let
\[
 P(z)=z^r+\sum_{j=0}^{r-1}c_jz^j,\qquad
 Q(z)=z^r+\sum_{j=0}^{r-1}\widehat c_jz^j.
\]
Choose $0<h<\min\{\gamma/2,m^2/2\}$. If
\begin{equation}\label{eq:rouche}
 \sum_{j=0}^{r-1}|\widehat c_j-c_j|(A^2+h)^j
           <h(\gamma-h)^{r-1},
\end{equation}
then $Q$ has exactly one root, counting multiplicity, in each disk
$|z-x_i|<h$. If its coefficients are real, all these roots are real
and positive. The corresponding half-lengths are within $h/m$ of the
true ones after matching.
\end{proposition}

\begin{proof}
On $|z-x_i|=h$, one factor of $P$ has modulus $h$ and all others have
modulus at least $\gamma-h$. The left side of \eqref{eq:rouche} bounds
$|Q-P|$ on that circle. Rouch\'e's theorem gives one root of $Q$ in each
of the disjoint disks. Real coefficients force nonreal roots to occur
in conjugate pairs; a disk centered on the real axis containing just
one root cannot contain a nonreal pair. Its real root is positive since
$h<m^2/2$. If $\widehat x_i$ is that root, then
\[
 |\sqrt{\widehat x_i}-\sqrt{x_i}|
 =\frac{|\widehat x_i-x_i|}{\sqrt{\widehat x_i}+\sqrt{x_i}}
 <\frac hm.
\]
\end{proof}

These are standard finite-dimensional mechanisms, included as a contrast
rather than as new root-perturbation theory. Their constants deteriorate
when nodes collide or their number grows. Both problems are present in
the bad Chebyshev blocks.

\subsection{What a certified computation would need}

A defensible finite-factor recovery procedure first specifies the number
of unknown factors, a lower bound on their lengths, a separation
condition, and a quantitative tail model. It then converts certified
moment intervals to cumulant intervals, subtracts the tail power-sum
intervals, applies Newton recursion with outward error bounds, and proves
root localization by a condition such as \eqref{eq:rouche}. Merely finding
numerical roots of a polynomial is not a certificate.

A useful elementary tail estimate is available without a detailed tail
formula. If its largest half-length is at most $\eta$ and its total
half-length is at most $L$, then
\begin{equation}\label{eq:tail-noise}
 \sum_{j\text{ in tail}}a_j^{2k}\le L\eta^{2k-1}.
\end{equation}
The proof factors out $\eta^{2k-1}$ and sums $a_j$. This can be treated
as an additional nonnegative uncertainty interval. It is not legitimate
to discard the tail because it contributes little variance: high-order
root recovery can remain sensitive to its structured contribution.

There is also a distinction between a fixed number of large factors and
the whole ordered sequence. The witnessing index in \cref{lem:gap} tends
to infinity. Consequently our lower bound does not, by itself, rule out
much better rates for estimating the first fixed $r$ half-lengths under
suitable tail assumptions.

\section{Reproducible verification and proof boundaries}
\label{sec:verification}

\subsection{What the included program actually verifies}

The package contains \texttt{code/verify.py}, requiring only Python,
SymPy, and mpmath. It uses integer Chebyshev recurrences and rational
Newton identities, avoiding numerical root finding for all algebraic
claims. A completed run records:
\begin{center}
\small
\begin{tabular}{@{}p{0.63\textwidth}r@{}}
\toprule
Check & Count\\
\midrule
Exact matching root power sums, before and after translation, $2\le n\le80$ & 6,320\\
Exact first unmatched power sums in both coordinates & 158\\
Exact Bernoulli cumulant defects, $2\le n\le40$ & 39\\
Exact Jacobian rank cases and augmented determinant cases, $1\le M\le8$ & 8 each\\
Explicit analytic curve invariants and fourth-cumulant derivative & passed\\
90-digit Chebyshev root residual diagnostics, $2\le n\le100$ & 10,098\\
90-digit middle-coordinate gap diagnostics & 99\\
90-digit small-frequency inequality diagnostics & 360\\
\bottomrule
\end{tabular}
\end{center}
The output files record software versions and the distinction between exact
checks and floating-point diagnostics. They deliberately do not describe
90-digit arithmetic as interval certification.

The finite tests do not prove an assertion for all $n$. Those proofs are
\cref{lem:cheby-powers,lem:gap,thm:tv}; the program is a reproducibility
and regression aid. In particular, the very small TV bounds are evaluated
from a proved formula rather than estimated by numerical quadrature of a
near-zero difference.

\subsection{Numerical scale of the construction}

\begin{table}[htbp]
\centering
\small
\begin{tabular}{@{}rrrr@{}}
\toprule
$n$ & Witness $\Delta_n$ & Fixed-variance witness $c_n\Delta_n$
    & Proved TV upper bound\\
\midrule
% The following rows are generated from data/verification.json and embedded
% so that this TeX file compiles without additional data files.
% Generated by code/verify.py; TV entries are analytic upper bounds.
16 & $5.781\times 10^{-4}$ & $5.204\times 10^{-4}$ & $1$ \\
32 & $1.462\times 10^{-4}$ & $1.384\times 10^{-4}$ & $7.755\times 10^{-2}$ \\
64 & $3.674\times 10^{-5}$ & $3.571\times 10^{-5}$ & $8.054\times 10^{-6}$ \\
128 & $9.206\times 10^{-6}$ & $9.074\times 10^{-6}$ & $4.343\times 10^{-14}$ \\
256 & $2.304\times 10^{-6}$ & $2.287\times 10^{-6}$ & $6.317\times 10^{-31}$ \\

\bottomrule
\end{tabular}
\caption{Witnessing coordinate gaps, not claimed exact $d_\infty$ distances.
The gap columns are rounded evaluations of explicit expressions; the last
column is the rounded value of the analytic upper bound $\eps_n$, not an
estimate of the actual total-variation distance.}
\label{tab:numeric}
\end{table}

For instance, the degree-$256$ fixed-variance pair has a displayed
coordinate separation about $2.29\times10^{-6}$ while the theorem bounds
its total variation by about $6.32\times10^{-31}$. This does not mean
that the actual total variation has been computed to that accuracy. It
means that an exact analytic estimate certifies an upper bound of that
size. The constants were chosen for a transparent proof, not numerical
sharpness.

\subsection{A possible Lean formalization route}

The results do not inherit formal status from the repository's existing
Fabius theory. A future formalization can nevertheless be decomposed into
clear interfaces. The first layer is finite algebra: Chebyshev recurrence,
the level-set root description, Newton identities, translated power sums,
and the middle-coordinate separation. These lemmas can be developed
without probability or infinite products.

The second layer represents the finite perturbations as product-measure
random variables and proves their moment and cumulant identities. The
third layer isolates the Fourier estimate as a statement about two finite
products, with separate small- and large-frequency lemmas and a
Plancherel-to-$L^1$ conversion. Only then is the common Fabius factor
introduced through a theorem interface giving its density, support,
variance, and derivative norms. The final layer is purely metric and
statistical: convolution contraction, dilation, local moduli, and the
two-point bound.

This dependency order keeps a large imported Fabius development out of
the finite algebraic core. It also makes the actual trust boundary
visible. No Lean source or uncompiled skeleton is supplied here, and no
machine-checking claim is made for the article.

\section{Further research questions}
\label{sec:questions}

\begin{question}[The sharp local inverse modulus]
Determine the asymptotic order of $\omega_r(\eps)$ in
\eqref{eq:modulus}. Does a logarithmic upper bound hold under the present
support, variance, and derivative constraints, and what exponent is
optimal? Our result proves only
$\omega_r(\eps)\gtrsim(\log(1/\eps))^{-2}$. Compactness gives convergence
to zero without a quantitative formula. A useful next step is to combine
truncated power-sum inversion with a uniform bound on the number of
coefficients above a chosen threshold.
\end{question}

\begin{question}[Pointwise recovery of the exact dyadic spectrum]
Decide whether a power-law estimate relative to the single reference
parameter $a^0$ can hold on a natural restricted class. The dyadic
spectrum has a strong separation pattern that the inserted Chebyshev
blocks destroy. A proof or counterexample must compare one law directly
to $\mu_0$, not merely compare two laws that separately approach it;
\cref{rem:pointwise} explains why the present argument is insufficient.
\end{question}
% ed. (2026-09-29): answered at the dyadic reference by a later package.
\begin{ednote}
The later repository article
\path{inverse-and-sampling/Anchored_Dyadic_Recovery_Uniform_Spectrum/} (filed
2026-09-29, unreviewed) answers this question negatively on the class $\K$
of \cref{def:K}, and also on $\A_L$ for $L>1$: the anchored modulus
$\sup\{d_\infty(a,a^0):a\in\K,\ \TV(\mu_a,\mu_0)\le\eps\}$ has the exact
scale $\exp[-\Theta(\sqrt{\log(1/\eps)})]$, and so does its analogue for
the Kolmogorov distance. So no positive H\"older exponent is possible, yet
the anchored modulus is eventually smaller than every negative power of
$\log(1/\eps)$, in particular than the lower bound
\eqref{eq:modulus-lower} for the pairwise modulus. Its leading constant is
bracketed, not determined.
\end{ednote}

\begin{question}[Exact support endpoints together with quantitative instability]
Construct explicit exponentially close pairs that retain
$\supp\mu=[-1,1]$ and variance $1/9$, approach $\mu_0$, and have a
quantified spectral separation. \Cref{thm:fibres} preserves the exact
support but does not give an exponentially small pairwise TV bound;
\cref{thm:main} gives that bound but only a common support containment.
A promising route is a small Chebyshev tail block with two compensating
large-factor adjustments, while controlling the errors those adjustments
introduce into low-order cumulants.
\end{question}

\begin{question}[Geometrically separated classes]
Assume $a_{j+1}\le q a_j$ for some fixed $q<1$, perhaps together with
bounds on the leading scale. Determine whether the inverse has a
substantially better modulus. The present construction is excluded,
since many factors are packed into one scale interval. Exact
identifiability by successive largest-power limits suggests an approach,
but converting it into certified finite-noise bounds requires controlling
every deflation error.
\end{question}
% ed. (2026-09-29): answered at the dyadic reference by a later package.
\begin{ednote}
At the dyadic reference only, the later article
\path{inverse-and-sampling/Anchored_Dyadic_Recovery_Uniform_Spectrum/} shows that
separation does not restore a power law: for every fixed $q\in(1/2,1)$ the
anchored modulus on $\{a\in\K:a_{j+1}\le qa_j\ \text{for all}\ j\}$ still
has the scale $\exp[-\Theta(\sqrt{\log(1/\eps)})]$. (The dyadic spectrum
itself needs $q\ge1/2$.) The pairwise modulus on separated classes and the
boundary case $q=1/2$ remain open.
\end{ednote}

\begin{question}[Fixed leading factors versus the complete sequence]
For fixed $r$, find sharp rates for the loss
$\max_{1\le j\le r}|\widehat a_j-a_j|$, with explicit assumptions on
the remaining tail. The middle-coordinate witness here occurs at a
rank growing with $n$, so it does not settle that finite-prefix question.
It is important to distinguish a separated finite-dimensional parameter
from an unknown infinite tail that can imitate its cumulants.
\end{question}
% ed. (2026-09-29): answered at the dyadic reference by a later package.
\begin{ednote}
Answered at the dyadic reference by the later article
\path{inverse-and-sampling/Anchored_Dyadic_Recovery_Uniform_Spectrum/}: for every $L\ge1$
and every fixed $r$, the prefix $(a_1,\dots,a_r)$ is locally
Lipschitz-stable at $a^0$ in Kolmogorov distance (hence in total variation)
over $\A_L$, with no assumption on the tail beyond the $\ell^1$ bound, and
the Lipschitz constant $C_{L,r}$ satisfies
$\log C_{L,r}\le(\log2)r^2+O_L(r+1)$. The optimal growth in $r$, and
prefix rates away from the dyadic reference, are not determined there.
\end{ednote}

\begin{question}[Matching statistical upper bounds]
Construct a regularized estimator with a proved risk bound over a
specified infinite-factor class, and compare it to
\cref{thm:statistics}. Empirical high moments become unreliable quickly;
a characteristic-function estimator with low-frequency regularization
may be more natural. Any proposed upper rate must account for both
moment cancellation and the number of unresolved small factors.
\end{question}
% ed. (2026-09-29): an anchored testing version is given by a later package.
\begin{ednote}
The later article
\path{inverse-and-sampling/Anchored_Dyadic_Recovery_Uniform_Spectrum/} gives only an
anchored version: testing $a=a^0$ against alternatives with
$d_\infty(a,a^0)\ge\delta$ in the geometrically separated subclasses of
$\K$ needs $\exp[\Theta(\log^2(1/\delta))]$ samples, and honest confidence
sets over $\A_L$ contract at $a^0$. A global estimator with a proved risk
bound over an infinite-factor class, as asked here, remains open.
\end{ednote}

\begin{question}[Gaussian components and square-summable spectra]
Extend the quantitative question to the larger generalized-uniform and
Gaussian-convolution parameter spaces studied by Billey--Swanson.
Their topology records the possibility that variance escapes into a
Gaussian component \cite{BilleySwanson}. Without an $\ell^1$ support
bound, new scale choices are possible, so the exponent and even the most
appropriate parameter metric may change. Establishing uniqueness alone
would repeat known work; a useful extension must supply quantitative
stability or instability information.
\end{question}
% ed. (2026-09-29): answered for the Gaussian-variance functional by a later
% package of the repository.
\begin{ednote}
For the Gaussian-variance functional, the later repository article
\path{inverse-and-sampling/Gaussian_Dust_Christoffel_Recovery_Uniform_Factors/} (filed 2026-09-29, unreviewed) answers this question in a
model with a known smooth compactly supported background: over
square-summable spectra with total latent variance at most $V$, the exact
minimax risk for the Gaussian variance is $V/2$ in absolute error at every
sample size, and a compact class admits uniformly consistent estimation
exactly when its squared-scale tails vanish uniformly. Quantitative
stability of the scales themselves in this larger space is not treated
there.
\end{ednote}

\begin{question}[Exact certificates for numerical inversion]
Implement a validated finite-prefix recovery algorithm based on
\cref{prop:largest,prop:rouche} and \eqref{eq:tail-noise}, then formalize
its algebraic certificate checker. The key engineering target is an
explicit statement of the measurement error model and tail assumptions,
not simply a high-precision reconstruction on selected examples. This
would turn the negative theorem into a practical guide for what a
reliable ProveIt inverse API must require.
\end{question}

\section{Conclusion}

The complete law of a summable uniform convolution uniquely encodes its
ordered factor lengths. That qualitative fact coexists with two strong
obstructions. A prescribed finite list of moments leaves nontrivial
analytic families even when the support is fixed exactly. More
quantitatively, a growing block of Chebyshev-derived positive factors can
hide polynomial parameter differences behind exponentially small errors
in the full probability law.

The construction retains fixed variance, compact support bounds, and
uniform control of every density derivative. Its instability therefore
lies in the factorization, not merely in rough densities or uncontrolled
scaling. The explicit proofs and algebraic checks support independent
review and eventual formalization. Sharp moduli, pointwise stability at
the dyadic law, and matching statistical algorithms remain the central
next questions.

\clearpage
\appendix
\section{A compact audit of the main constants}
\label{app:constants}

The following chain contains all numerical losses used in the main
construction:
\begin{align*}
 &\frac1{4n}\le\eta_n\le\frac1{2n},\qquad
   |y_{j_n}^+-y_{j_n}^-|\ge\frac1{3n},\\
 &\Delta_n\ge\frac{\eta_n}{12n}\ge\frac1{48n^2},
   \qquad\frac12<c_n\le1,\\
 &\int_{|t|\le1}(|D|^2+|D'|^2)\dd t\le8320n^2q^{2n},\\
 &\int_{|t|\ge1}(|D|^2+|D'|^2)\dd t\le80n^2\rho^{4n-2},\\
 &q=4/\pi^2<1/2,\quad\rho=6/7,\quad
   q^{2n}\le\rho^{4n-2},\\
 &\TV\le\frac{\sqrt{8400}}{2\sqrt2}n\rho^{2n-1}
                       <40n\rho^{2n-1}.
\end{align*}
The normalization changes the gap lower bound by at most a factor of two
and leaves the TV bound unchanged. The statistical argument loses a
further factor of eight: four from the testing reduction and two from
the product-TV bound $1-\TV\ge1/2$. Thus its denominator is
$96\cdot8=768$.

The local-modulus denominator is $96\cdot11^2=11616$. Its threshold in
$\eps$ depends on the prescribed neighborhood radius only through the
explicit convergence estimate \eqref{eq:main-near}. None of these
constants is claimed to be optimal.

\begin{thebibliography}{9}
\small\raggedright

\bibitem{ProveIt}
Vladimir Reshetnikov and contributors.
\emph{ProveIt}, especially \texttt{Analysis/FabiusFunction/README.md}
and its Fabius/Rvachev documentation.
Repository snapshot reference:
\texttt{\repoSHA}.
\url{https://github.com/VladimirReshetnikov/ProveIt}.
\href{https://github.com/VladimirReshetnikov/ProveIt/blob/74f7f5bdba1aa728b340509701a0fa4e45e8dbf3/Analysis/FabiusFunction/README.md}{Project overview at the inspected snapshot}.
Targeted documentation review, not a complete repository proof audit.

\bibitem{AriasDefinition}
Juan Arias de Reyna.
\emph{An infinitely differentiable function with compact support:
Definition and properties}.
arXiv:1702.05442, 2017. English translation of
``Definici\'on y estudio de una funci\'on indefinidamente diferenciable de
soporte compacto,'' \emph{Rev. Real Acad. Ciencias} \textbf{76} (1982),
21--38.
\url{https://arxiv.org/abs/1702.05442}.

\bibitem{AriasArithmetic}
Juan Arias de Reyna.
\emph{Arithmetic of the Fabius function}.
arXiv:1702.06487, version 3.
\url{https://arxiv.org/abs/1702.06487}.
Used as repository context rather than as an input to the instability proof.

\bibitem{BilleySwanson}
Sara C. Billey and Joshua P. Swanson.
\emph{The metric space of limit laws for $q$-hook formulas}.
\emph{Combinatorial Theory} \textbf{2}(2), 2022.
DOI: \href{https://doi.org/10.5070/C62257868}{10.5070/C62257868}.
The theorem numbering cited here is that of
arXiv:2010.12701v2, 2023, especially Lemma~3.11,
Theorem~3.13, and Theorem~1.15.
\url{https://arxiv.org/abs/2010.12701v2}.

\bibitem{GerthEtAl}
Daniel Gerth, Bernd Hofmann, Christopher Hofmann, and Stefan Kindermann.
\emph{The Hausdorff Moment Problem in the light of ill-posedness of type I}.
arXiv:2104.06029v2, 2021.
\url{https://arxiv.org/abs/2104.06029v2}.
Cited for related moment-inversion context; no theorem of that paper is
used as a substitute for the explicit uniform-factor construction.

\end{thebibliography}
\end{document}
